\bookchapter{الفصل 2: رتّب غرفتك الرقمية — إدارة الملفات}

\chapterintro{الهاتف مع الاستخدام يصبح فوضوياً، والعادات الجيدة تبدأ من إدارة الملفات. كما أن غرفتك تحتاج إلى ترتيب دوري، كذلك مساحتك الرقمية تحتاج إلى عناية. هاتف منظم لا يسهّل العثور على الأشياء فحسب، بل يعمل بسلاسة أكبر أيضاً.}

\booksection{أولاً: نظام التخزين ثلاثي الطبقات: مستودعك الرقمي}

\begin{enumerate}
  \item \textbf{تخزين الجسم}: المنطقة الأساسية، سرعة عالية، النظام والتطبيقات كلها هنا. هذه أثمن مساحة في هاتفك، يجب استخدامها بحكمة.
  \item \textbf{بطاقة SD (إذا كانت مدعومة)}: خزانة ملحقات، أبطأ، مناسبة للأفلام والصور والملفات الكبيرة. ملاحظة: العديد من الهواتف الرئيسية لم تعد تدعم بطاقات SD.
  \item \textbf{التخزين السحابي}: خزنة + قرص شبكي للمزامنة. يمنع الفقد، ويمكنه المزامنة بين أجهزة متعددة. iCloud و Baidu Netdisk و OneDrive كلها خيارات شائعة.
\end{enumerate}

\booksection{ثانياً: تعرّف على “مدير الملفات”}

Android عادةً لديه تطبيق [الملفات] أو [إدارة الملفات]، iOS لديه تطبيق [الملفات] المدمج. إنه “مستكشف الموارد” لعرض جميع ملفات هاتفك.

\booksubsection{معاني المجلدات الشائعة:}

\begin{itemize}
  \item \textbf{DCIM}: الصور والفيديوهات التي التقطتها الكاميرا.
  \item \textbf{Pictures}: صور أخرى، مثل لقطات الشاشة والصور المحفوظة.
  \item \textbf{Download}: الملفات التي نزّلها المتصفح.
  \item \textbf{Music}: ملفات الموسيقى.
  \item \textbf{Movies}: ملفات الفيديو.
  \item \textbf{Documents}: ملفات المستندات.
  \item \textbf{Android}: مجلد نظام Android، بداخله بيانات التطبيقات، لا تحذفه يدوياً عادةً.
\end{itemize}

\booksection{ثالثاً: بنية نظام الملفات: فهم أعمق}

Android مبني على Linux، ونظام الملفات له بنية هرمية معينة. رغم أن المستخدم العادي لا يحتاج إلى فهم عميق، إلا أن معرفة بعض المفاهيم الأساسية مفيدة.

\booksubsection{أقسام التخزين:}

\begin{itemize}
  \item \textbf{/system}: قسم النظام، يخزن ملفات نظام التشغيل. المستخدم العادي لا يملك صلاحية الوصول، يحتاج إلى صلاحية root.
  \item \textbf{/data}: قسم البيانات، يخزن بيانات المستخدم وبيانات التطبيقات. هذه هي المساحة التي تستخدمها يومياً.
  \item \textbf{/sdcard أو /storage/emulated/0}: مساحة التخزين التي يمكن للمستخدم الوصول إليها، صورك وملفات التنزيل كلها هنا.
\end{itemize}

\booksubsection{تخزين بيانات التطبيقات:}

\begin{itemize}
  \item \textbf{التخزين الداخلي}: المساحة الخاصة بالتطبيق، لا يمكن للتطبيقات الأخرى الوصول إليها. تُحذف عند إزالة التطبيق.
  \item \textbf{التخزين الخارجي}: الملفات التي يخزنها التطبيق في المنطقة التي يمكن للمستخدم الوصول إليها، مثل الصور ومحتوى التنزيل. لا تُحذف بالضرورة عند إزالة التطبيق.
\end{itemize}

\booksection{رابعاً: النسخ الاحتياطي للبيانات واستعادتها: تأمين بياناتك}

البيانات لا تقدّر بثمن، والنسخ الاحتياطي عادة يجب أن تتبناها.

\booksubsection{استراتيجية النسخ الاحتياطي:}

\begin{enumerate}
  \item \textbf{النسخ السحابي}: تلقائي، مريح، لكن يتطلب شبكة ومساحة سحابية. iCloud و Huawei Cloud و Xiaomi Cloud وغيرها تدعمه.
  \item \textbf{النسخ المحلي}: الاتصال بالحاسوب بكابل بيانات، النسخ بأداة رسمية. سريع، مساحة كبيرة، لكن يتطلب تشغيلاً يدوياً.
  \item \textbf{النسخ المزدوج}: نسخ سحابي + نسخ محلي، الأكثر أماناً.
\end{enumerate}

\booksubsection{ماذا تنسخ احتياطياً؟}

\begin{itemize}
  \item \textbf{الصور والفيديوهات}: الأكثر احتلالاً للمساحة، والأثمن.
  \item \textbf{جهات الاتصال}: الأفضل تخزينها في حساب سحابي.
  \item \textbf{سجلات محادثة WeChat}: استخدم وظيفة النسخ الاحتياطي المدمجة في WeChat.
  \item \textbf{المستندات المهمة}: خزّنها في قرص سحابي أو على حاسوب.
\end{itemize}

\booksection{خامساً: أمان الملفات: حماية خصوصيتك}

\begin{itemize}
  \item \textbf{تشفير الملفات الحساسة}: العديد من الهواتف لديها وظيفة "المجلد الآمن" أو "المساحة الخاصة"، يمكن وضع الملفات الحساسة فيها، وتتطلب كلمة مرور أو بصمة للوصول.
  \item \textbf{الحذف لا يعني التدمير}: الحذف العادي فقط يعلّم المساحة قابلة للكتابة فوقها، البيانات لا تزال موجودة. للحذف التام، استخدم أدوات مسح آمن مخصصة، أو اكتب فوقها عدة مرات.
  \item \textbf{الأمان السحابي}: اختر مزود خدمة سحابية ذا سمعة جيدة، فعّل المصادقة الثنائية، وراجع بانتظام الأجهزة المسجلة.
\end{itemize}

\booksection{سادساً: المساحة لا تكفي؟ قم بتنظيف كبير}

\begin{itemize}
  \item \textbf{تنظيف ذاكرة التخزين المؤقت}: الذاكرة المؤقتة للتطبيقات، تنظيفها يفتح مساحة. في الإعدادات > التخزين يمكن التنظيف.
  \item \textbf{البحث عن الملفات الكبيرة}: استخدم مدير الملفات للترتيب حسب الحجم، اعثر على "المذنبين" الذين يحتلون المساحة.
  \item \textbf{حذف البقايا}: بعد إزالة التطبيقات، أحياناً تترك ملفات غير مفيدة، احذفها يدوياً.
  \item \textbf{تنظيف الألبوم}: احذف الصور المكررة، الصور الضبابية، لقطات الشاشة غير الضرورية.
  \item \textbf{تنظيف WeChat}: WeChat مستهلك كبير للمساحة، تنظيف سجلات المحادثة وذاكرة التخزين المؤقت دورياً ضروري جداً.
\end{itemize}

\booksection{سابعاً: تحليل شامل لتنسيقات الملفات الشائعة}

امتدادات الملفات المختلفة في الهاتف مربكة، فهمها يساعدك على إدارة الملفات بشكل أفضل:

\begin{center}
\begin{tabularx}{\linewidth}{|X|X|X|X|}
\hline
\textbf{التنسيق} & \textbf{النوع} & \textbf{الخصائص} & \textbf{طريقة الفتح الموصى بها} \tabularnewline
\hline
JPG/JPEG & صورة & ضغط مع فقد، حجم صغير، أفضل توافق. التنسيق الافتراضي لصور الهاتف. & الألبوم/أي عارض صور \tabularnewline
PNG & صورة & ضغط بدون فقد، يدعم الخلفية الشفافة. شائع للقطات الشاشة والأيقونات. حجم أكبر من JPG. & الألبوم/أي عارض صور \tabularnewline
MP3 & صوت & ضغط مع فقد، حجم صغير، توافق ممتاز. تنسيق الصوت الأكثر شيوعاً. & أي مشغل موسيقى \tabularnewline
FLAC & صوت & ضغط بدون فقد، جودة صوت مطابقة لـ CD، حجم أكبر من MP3 بـ 3-5 مرات. الخيار الأول للهواة. & مشغل موسيقى/تطبيق تشغيل احترافي \tabularnewline
WAV & صوت & صوت خام بدون ضغط، أفضل جودة لكن حجم ضخم. شائع للتسجيلات الاحترافية. & مشغل موسيقى \tabularnewline
M4A & صوت & ترميز AAC، جودة صوت أفضل من MP3 عند نفس معدل البت. شائع في نظام Apple البيئي. & Apple Music/مشغل موسيقى \tabularnewline
MP4 & فيديو & تنسيق تغليف الفيديو الأكثر شيوعاً، توافق ممتاز. التنسيق الافتراضي لفيديو الهاتف. & الألبوم/أي مشغل فيديو \tabularnewline
MKV & فيديو & تنسيق تغليف مفتوح المصدر، يدعم مسارات صوت متعددة، ترجمات متعددة. شائع لموارد الأفلام. & VLC/nPlayer \tabularnewline
DOC/DOCX & مستند & مستند Word. DOCX تنسيق مضغوط، حجم أصغر. DOC تنسيق قديم. & WPS/Office \tabularnewline
PPT/PPTX & عرض تقديمي & عرض تقديمي PowerPoint. PPTX تنسيق مضغوط. & WPS/Office \tabularnewline
XLS/XLSX & جدول & جدول Excel الإلكتروني. XLSX تنسيق مضغوط. & WPS/Office \tabularnewline
PDF & مستند & تنسيق عبر الأنظمة، تخطيط ثابت غير قابل للتحرير. شائع للعقود والأطروحات والكتب الإلكترونية. & Adobe Reader/WPS \tabularnewline
ZIP & مضغوط & تنسيق الضغط الأكثر شيوعاً، توافق عبر الأنظمة جيد. & مدمج في النظام/تطبيق فك ضغط \tabularnewline
RAR & مضغوط & نسبة ضغط أعلى من ZIP، لكن يتطلب برنامجاً مخصصاً لفك الضغط. & تطبيق فك ضغط \tabularnewline
ISO & صورة & ملف صورة قرص، ينسخ محتوى القرص بالكامل. شائع لحزم تثبيت النظام. & تطبيق فك ضغط/محرك أقراص افتراضي \tabularnewline
HTML & صفحة ويب & ملف صفحة ويب، يفتح بالمتصفح. الصفحات المحفوظة شائعة بهذا التنسيق. & المتصفح \tabularnewline
\hline
\end{tabularx}
\end{center}

\begin{tcbpractice}{تمرّن: رتّب مجلد "التنزيلات"}
افتح مدير الملفات > ادخل إلى مجلد “التنزيلات”. انظر ماذا يوجد بداخله؟ احذف غير المفيد، وصنّف المفيد (مثلاً أنشئ مجلدات “مواد العمل” و“رموز تعبيرية” وضعها فيها). بالمناسبة، انظر إلى مساحة التخزين لمعرفة أنواع الملفات التي تحتل المساحة أكثر.

\end{tcbpractice}

\booksection{ثامناً: إدارة الصور ونظام الملفات: اجعل الأصول منظمة}

في القسم السابق تعلمنا “التعرف على المجلدات”، لكن التعرف وحده لا يكفي. هذا القسم عن\textbf{كيفية بناء مواصفات إدارة قابلة للتنفيذ طويل الأمد بشكل استباقي}، بحيث تجد ملفاتك دائماً ولا تفقدها ولا يتم التجسس عليها.

\booksubsection{تصنيف الملفات وقواعد التسمية}

جذر الفوضى هو غياب القواعد. بناء القواعد بسيط في الواقع، المفتاح هو\textbf{الاستمرارية}.

\begin{itemize}
  \item \textbf{طريقة تسمية "الوقت + الحدث"}: اسم المجلد يحمل طابعاً زمنياً ووصفاً للحدث، يسهّل التذكر والبحث لاحقاً. مثلاً\texttt{202606\_تذاكر\_السفر}و\texttt{ملخص\_2025\_السنوي}و\texttt{2024\_عيد\_ميلاد\_الطفل}. هكذا يمكنك رؤية الخط الزمني، ومعرفة المحتوى بنظرة واحدة.
  \item \textbf{نظام تصنيف ثلاثي المستويات}:
  \begin{itemize}
    \item \textbf{المستوى الأول}: حسب نوع الملف (مستند / صورة / فيديو / صوت)
    \item \textbf{المستوى الثاني}: حسب سيناريو الاستخدام (عمل / حياة / تعلم)
    \item \textbf{المستوى الثالث}: حسب حدث أو مشروع محدد (مثل "2026Q2\_مشروع\_A")
  \end{itemize}
  \item \textbf{خطة التنظيف الشهرية}: حدد يوماً ثابتاً كل شهر (يُنصح يوم الراتب أو اليوم الأول من الشهر)، للتنظيم المركز لمجلد "التنزيلات":\textbf{أرشفة المفيد، حذف غير المفيد، نقل الثمين}. 12 تنظيماً صغيراً في السنة أسهل بكثير من تنظيف كبير بعد سنة كاملة.
\end{itemize}

\begin{tcbpractice}{قائمة مهام هذا الشهر}
اختر عطلة نهاية أسبوع فارغة، واعمل بالترتيب التالي:

\begin{enumerate}
  \item أنشئ في مدير الملفات مجلد عمل لهذا الشهر بشكل\texttt{2026Q2\_مشروع\_A}؛
  \item راجع جميع ملفات هذا الأسبوع في مجلد "التنزيلات"، احتفظ بالمفيد، أرشفة، احذف غير المفيد؛
  \item أنشئ لصور العائلة لهذا العام مجلدات سنوية مثل\texttt{عائلة\_2025}و\texttt{عائلة\_2026}.
\end{enumerate}

\end{tcbpractice}

\booksubsection{استراتيجية معالجة الملفات الزائدة}

الهاتف مع الاستخدام “يزداد وزناً”، غالباً بسبب الملفات الزائدة.

\begin{itemize}
  \item \textbf{التنظيف التلقائي خلال 72 ساعة}: الملفات المؤقتة التي تستقبلها تطبيقات التواصل الاجتماعي مثل WeChat و QQ (رموز تعبيرية، صور الدردشة، فيديوهات صغيرة)، معظمها يفقد فائدته بعد مشاهدة واحدة. في [WeChat] > [الإعدادات] > [عام] > [مساحة التخزين] فعّل "التنظيف التلقائي"، اضبط الذاكرة المؤقتة منتهية الصلاحية على 72 ساعة. الهاتف يصبح أنظف بكثير، دون فقدان سجلات مهمة.
  \item \textbf{التنظيف الموجه لمساحة التخزين}: Android [الإعدادات] > [التخزين] > [تنظيف وتسريع]، iOS [الإعدادات] > [عام] > [مساحة تخزين iPhone]، سيصنف النظام "ذاكرة التخزين المؤقت / الملفات المؤقتة / التطبيقات القابلة للإزالة / الملفات الكبيرة". ركّز على تنظيف: (1) ملفات الذاكرة المؤقتة؛ (2) حزم التثبيت منتهية الصلاحية (التي تترك بعد ترقية التطبيقات)؛ (3) الملفات المكررة (صور بنفس الاسم والحجم).
  \item \textbf{قواعد التصنيف التلقائي}: معظم تطبيقات الألبوم للعلامات التجارية تدعم التصنيف الذكي: الأشخاص، المناظر الطبيعية، المستندات، لقطات الشاشة. بعد التفعيل، الصور الجماعية الجديدة تُصنف تلقائياً في "الأشخاص > Zhang San"، والعقود الممسوحة تُصنف تلقائياً في "المستندات".\textbf{بالتعاون مع التنظيم الشهري، يمكن تقليل احتمال "عدم العثور على الأشياء" بشكل كبير}.
\end{itemize}

\booksubsection{مواصفات معالجة الصور الآمنة}

الصور هي أكثر أنواع البيانات “حساسية” في الهاتف — صورة عابرة يمكن أن تكشف موقعك وجهازك ووقتك. ننصح بـ\textbf{بناء سير عمل “حماية مزدوجة”}:

\begin{tcbwarn}{الإجراء القياسي لإرسال صور الهوية / لقطات الشاشة الحساسة}
\begin{enumerate}
  \item \textbf{طمس المعلومات المرئية}: استخدم أداة مدمجة أو طرف ثالث (مثل Meitu XiuXiu، Snapseed)\textbf{لطمس أرقام الهوية الوسطى، رقم اللوحة، رقم الباب، الاسم، المبلغ}وغيرها من الحقول الرئيسية؛
  \item \textbf{إزالة البيانات الوصفية}: قبل الإرسال، اضغط على "مشاركة" أو "تصدير" في الألبوم، سيسأل النظام عما إذا كنت تريد الاحتفاظ بمعلومات الموقع، اختر\textbf{"عدم الاحتفاظ"}. هكذا لا يتم إرسال طراز جهاز التصوير، إحداثيات GPS، وقت التصوير مع الصورة؛
  \item \textbf{علامة مائية مؤقتة}(اختياري): أضف علامة مائية نصية مثل "للاستخدام في معاملة XX فقط" على لقطة الشاشة الحساسة، لمنع إساءة الاستخدام.
\end{enumerate}

\end{tcbwarn}

طرق شائعة لتنظيف البيانات الوصفية (EXIF) (عمليات أساسية، للحماية المتقدمة راجع الفصل 6 “إزالة معلومات الصور الحساسة”):

\begin{itemize}
  \item \textbf{مدمج في النظام}: iOS [الإعدادات] > [الخصوصية والأمان] > [الكاميرا]، يمكن إيقاف "الموقع الدقيق"؛ عند المشاركة اختر "الصور الحالية فقط"\textbf{سيتم إزالة معظم EXIF تلقائياً}.
  \item \textbf{أندرويد}: معظم تطبيقات الألبوم للعلامات التجارية تطلب عند "المشاركة" "هل تريد إزالة معلومات الموقع"، حدد ذلك.
  \item \textbf{أدوات احترافية}: ExifTool (على الحاسوب)، Metapho (iOS)، Photo Metadata Remover (Android) يمكنها إزالة جميع البيانات الوصفية بالكامل.
\end{itemize}

بالنسبة للصور الحساسة التي تحتاج\textbf{حفظاً طويل الأمد}(بطاقة الهوية، البطاقة البنكية، نسخ مسحية لسجل العائلة، جواز السفر)، يُنصح بتخزينها في\textbf{ألبوم مشفّر}:

\begin{itemize}
  \item iPhone: [الصور] > حدد الصورة > إخفاء (يمكن عرضها فقط بعد التحقق بـ Face ID)؛
  \item Huawei/Xiaomi: [الألبوم] > اضغط مطولاً على الصورة > "نقل إلى الخزنة"؛
  \item تطبيقات تشفير طرف ثالث:\textbf{ألبوم 360 الآمن}و\textbf{مدير ألبوم Tencent (ألبوم الخصوصية)}و\textbf{KeepSafe}وغيرها، حماية بكلمة مرور منفصلة، إخفاء أفضل.
\end{itemize}
